\subsection*{Output Transformations}
\label{sec:output-transformations}

Transformations of the target values are a common technique to improve the performance of the surrogate model~\cite{jones1998}.
When the surrogate model is trained on the transformed target values, one would typically apply the inverse transformation to the predictive mean and standard deviation before passing them to the acquisition function, to get values on the original scale.
Alternatively, predictions can be kept on the transformed scale and the acquisition function can be defined directly in that space.
LogEI, for example, is explicitly designed to operate on log-transformed targets.

\paragraph{Log-Transformed Surrogate Models}

In many hyperparameter optimization scenarios, the distribution of the target variable (e.g., runtime or loss) is heavy right-tailed due to some very poorly performing configurations.
These configurations can correspond to timeouts or models that fail to train and are typically assigned very large loss or runtime values.
Log-transforming the target values compresses these extremes, makes the marginal distribution more symmetric and Gaussian-like, and can reduce heteroscedasticity.
However, since target values are required to be positive for log transformation, it may be possible to add an offset when negative objective values can be expected.
\texttt{mlr3mbo} therefore scales output values to $(0, 1)$ before applying the log transformation.

GP surrogate models tend to be more accurate in this setting when they are trained on $\log(y)$ instead of $y$~\cite{jones1998,hutter2009}, because the GP assumptions of Gaussian, homoscedastic residuals are better satisfied.
Similarly, random forests tend to find more balanced splits on log-transformed targets and can perform better on heavy-tailed objectives~\cite{lindauer2019}.

When LCB or PI are computed directly on the log-transformed surrogate, no inverse transformation of the predictions is required and the acquisition function is simply interpreted as operating in log-space.
If, instead, these acquisition functions are to be applied on the original scale $y$, the predictive distribution in log-space must be mapped back to $y$ (e.g., using the log-normal moment formulas).

\paragraph{LogEI Acquisition Function}

The LogEI acquisition function is the expected improvement on the original objective scale when the surrogate is trained on log-transformed target values~\cite{hutter2009}. This means it is
\begin{displaymath}
\alpha_{\mathrm{LogEI}}(\xv) = \mathbb{E} \left[ \max \left( f_{\min} - Y(\xv), 0 \right) \right],
\end{displaymath}
with $Y(\xv)$ interpreted as a log-normal distribution and $\ln Y(\xv) \sim \mathcal{N}(\mu(\xv), \sigma^{2}(\xv))$.
The LogEI acquisition function is then defined as
\begin{equation}
  \alpha_{\mathrm{LogEI}}(\xv)
  = f_{\min}\,\Phi\big(v(\xv)\big)
    - \exp\!\Big(\mu(\xv) + \tfrac{1}{2}\sigma^{2}(\xv)\Big)\,
      \Phi\big(v(\xv) - \sigma(\xv)\big),
\end{equation}
where
\begin{equation}
  v(\xv) = \frac{\ln f_{\min} - \mu(\xv)}{\sigma(\xv)}.
\end{equation}
Here, $\Phi$ denotes the cumulative distribution function of the standard normal distribution, and $f_{\min}$ is the best function value in the archive.
Thus, LogEI computes the expected improvement on the original scale under a log-normal model, which can better capture heavy right tails than a surrogate model defined directly on $y$.
Both the original SMAC and SMAC3 (in its SMAC4HPO facade) use LogEI as the default acquisition function whenever a log-transform of the performance values is employed~\cite{hutter2011,lindauer2022}. LogEI should not be confused with the logarithm of the EI, which is also used by various methods~\cite{watanabe2024}.

